%% kapa: exam=1 round=25 year=2014 period=1 form=A subject=economics subject_ko=경제원론 seq=35 number=75
%% source: Q-Net 감정평가사 기출문제 (cst00309) artlSeq=5101207 file=extracted/1008/5101207/2014년 제25회 감정평가사 1차시험 기출문제/1교시(A형).hwp
%% conversion: kordoc 4.24.51 -> tools/gen1.py (auto); review=proofread-sonnet; flags=none
\begin{question}{75}
화폐수요함수는 $\frac{ M } { P } =L(i, Y)$이고, $i$가 낮을수록, $Y$가 높을수록 화폐수요는 증가한다. 중앙은행이 내년부터 통화량 증가율을 높이기로 발표할 때 모든 개인들이 이 발표를 그대로 신뢰한다면, 금년도 이 경제에서 발생되는 현상으로 옳은 것은? (단, $P$는 물가, $M$은 통화량, $r$은 실질이자율, $\pi^{ e }$은 예상 물가상승률, 명목이자율 $i=r+ \pi^{ e }$, $Y$는 실질GDP이며, 물가는 항상 신축적으로 조정된다.)
\begin{choices}
  \item 아직 통화량 증가가 발생한 것은 아니므로 $\pi^{ e }$도 변하지 않을 것이다.
  \item 아직 통화량 증가가 발생한 것은 아니지만 현재 물가 $P$는 하락할 것이다.
  \item 아직 통화량 증가가 발생한 것은 아니지만 현재 물가 $P$는 상승할 것이다.
  \item $\pi^{ e }$가 상승할 것이므로 화폐수요가 증가할 것이다.
  \item $\pi^{ e }$가 하락할 것이므로 화폐공급이 증가할 것이다.
\end{choices}
\end{question}
\answer{3}
